
Hypothesis testability assessment in deep learning research currently relies on informal expert judgment or trial-and-error discovery. A researcher formulating a fairness hypothesis requiring labeled demographic annotations may invest weeks before discovering the target dataset lacks required metadata---a constraint violation detectable in minutes through formal verification. Existing dataset catalogs such as Papers With Code and HuggingFace Datasets Hub provide resource discovery but lack testability classification logic. Expert judgment systems achieve 80-85\% inter-rater reliability but validate predictions circularly against other expert opinions rather than experimental outcomes.

This work treats testability in constraint-driven contexts as a formal constraint-satisfiability problem: ∃ (Dataset D, Benchmark B, Metric M) such that hypothesis H can be expressed as a measurable intervention M(H\_intervention) - M(H\_baseline). Absence of any element renders the hypothesis untestable. This formalism enables verification via existence checking over a structured knowledge base extracted from catalog APIs.

We present a constraint-satisfiability verification system with three contributions: (1) Automated knowledge base construction from the HuggingFace Datasets Hub API achieves 84\% coverage of well-known deep learning datasets without manual curation. (2) Cross-domain confound pattern flagging achieves 93.33\% precision by transferring 15 documented patterns from literature across modalities via keyword-based matching. (3) Experimental ground truth validation---testing 20 system-classified ''testable'' hypotheses and measuring post-hoc p-values---demonstrates 90\% experimental success rate (binomial p = 0.0002 vs. 50\% random baseline), providing proof-of-concept non-circular evaluation for meta-research tools.

Formal verification achieves 0\% false positive rate through conservative (D,B,M) existence checking. Domain boundary detection correctly flags 100\% of out-of-scope hypotheses (10/10 novel modalities) before knowledge base lookup. Two testable hypotheses yielded null results (p = 0.679, p = 0.757)---these represent legitimate negative findings where (D,B,M) triples existed and no confounds were present, but interventions did not produce significant effects.
